\subsection{Justificativa}

\par Os processos industriais geram grande volume de dados, utilizados para acompanhar variáveis e interpretar a tendência do processo. Em geral, essa informação é organizada e apresentada ao operador por um software de supervisão, por meio de uma interface homem-máquina (IHM), que converte leituras de campo em telas, tendências e alarmes de forma legível e acionável \cite{ZANGHI2019_SCADA_CONCEITOS}. Com o aumento da capacidade de processamento dos computadores, o aprimoramento dos protocolos de comunicação e a maior integração com redes externas, inclusive a Internet, consolidou-se a supervisão remota e cresceu a exigência de visibilidade em tempo real sobre o que ocorre na planta industrial. Em setores como saneamento, transporte e energia, nos quais instalações costumam estar geograficamente espalhadas, ganham peso a telemetria e o controle supervisionado à distância. De modo esquemático, uma arquitetura típica de sistema SCADA integra os equipamentos de campo, como CLPs, aos módulos de monitoramento e controle e centraliza a visualização dos dados para o operador \cite{AHMED2011_EFFICIENT_WEB_SCADA}.

\par Apesar de consolidados, os sistemas supervisórios convencionais costumam exigir investimento elevado em licenças de software proprietário, hardware dedicado e infraestrutura local, o que os mantém fora do alcance de pequenas e médias empresas, para as quais a aquisição de uma licença raramente se justifica diante do porte da operação \cite{PHUYAL2020_COST_EFFICIENT_SCADA}. Nesse cenário, oferecer a supervisão como um serviço web, substitui esse investimento inicial por uma assinatura sob demanda, alinhada ao uso efetivo do sistema, e democratiza o acesso a uma tecnologia antes economicamente inviável para esse público. Além de reduzir custos, a supervisão via web amplia o alcance do monitoramento e do controle a instalações geograficamente dispersas e favorece a colaboração entre equipes em locais distintos, reduzindo deslocamentos presenciais e agilizando a resposta a ocorrências \cite{AHMED2011_EFFICIENT_WEB_SCADA}.

\par A computação em nuvem é um modelo em que recursos computacionais, como armazenamento, processamento, bancos de dados e aplicações, são ofertados como serviços remotos, flexíveis e escaláveis \cite{PEDROSA2011_COMPUTACAO_NUVEM}. Em geral, o usuário não precisa conhecer a localização física nem os detalhes da infraestrutura que sustenta esses serviços, o consumo costuma ser sob demanda, alinhado ao uso, o que reduz a necessidade de investimento direto em hardware e software no ambiente local \cite{HENRIQUES2010_CLOUD_COMPUTING_TCC}. Por sua vez, a conectividade, em especial por meio de redes sem fio, é essencial para que dispositivos e sistemas ciber-físicos mantenham troca de dados com a nuvem e com demais elementos da arquitetura distribuída \cite{STANKOVIC2014_IOT_RESEARCH_DIRECTIONS,MONOSTORI2016_CPS_MANUFACTURING}.